
\cleardoublepage
\phantomsection
\chapter*{\Huge\bfseries\hfill\textRL{مـلـخـص}}
\thispagestyle{noheader}
\addcontentsline{toc}{chapter}{\textRL{مـلـخـص}}



\vspace{0.1cm}





\vspace{0.1cm}
\begin{RLtext}

يُعد القمح من المحاصيل الزراعية الأساسية على مستوى العالم، إلا أن إنتاجه يواجه تهديدًا كبيرًا بسبب تنوع الأمراض التي تصيبه. تُعتبر طرق التشخيص التقليدية، المعتمدة على الفحص البصري اليدوي، بطيئة، وذات طابع ذاتي، وغالبًا ما تكون غير متاحة للمزارعين، خاصة في المناطق النائية. ومع ظهور الزراعة الذكية وتطور تقنيات التعلم العميق، أصبحت عملية التعرف التلقائي على الأمراض من خلال الصور حلاً واعدًا.

يهدف هذا المشروع إلى تصميم وتنفيذ نظام فعّال يعتمد على الشبكات العصبية الالتفافية \CNN للكشف عن أمراض القمح وتصنيفها باستخدام صور لمختلف أجزاء النبات مثل الأوراق والسيقان والسُنابل. تشمل منهجيتنا مراجعة شاملة للدراسات السابقة، وجمع البيانات ومعالجتها، وتدريب النماذج وتقييمها، ثم دمج الحل في منصة ويب تفاعلية موجهة للمزارعين. ولتعزيز فاعلية النظام، قمنا أيضًا بإدراج وحدة لاكتشاف الحشرات الضارة ومساعد ذكي يعمل عبر واجهة برمجة التطبيقات (API) لتوجيه المستخدمين.

نقترح إطارًا هرميًا للتصنيف يتكون من مرحلتين: في المرحلة الأولى، يتم دمج عدة هياكل CNN بشكل فردي مع خوارزمية آلة متجه الدعم (SVM) للتمييز بين العينات السليمة والمصابة، مشكلين بذلك مقاربة هجينة CNN-SVM. أما في المرحلة الثانية، فنستخدم بنية ResNet50 معززة بوحدة الانتباه الكفء للقنوات (ECA) للتعرف على تسع أمراض محددة تصيب القمح. أظهرت النتائج التجريبية أن المصنف الثنائي حقق دقة بلغت 99{،}60\%، بينما بلغ المصنف متعدد الأصناف دقة قدرها 99{،}83\%. وبلغت الدقة الإجمالية للنظام الكامل 99{،}48\%.

تُساهم الحلول المقترحة في تطوير الزراعة الذكية من خلال تمكين التشخيص المبكر والموثوق، مما يساعد في التدخل في الوقت المناسب وتحسين إدارة المحاصيل الزراعية.

\end{RLtext}

\vspace{0.1cm}
\hspace*{0mm}\rule{\linewidth}{1pt}
\vspace{0.1cm}
\begin{RLtext}
\noindent\textbf{الكلمات المفتاحية —} الزراعة الذكية، التعلُّم العميق، الشبكات العصبية الالتفافية، الرؤية الحاسوبية، اكتشاف أمراض القمح.
\end{RLtext}
\vspace{0.1cm}
\rule{\linewidth}{1pt}


